
% -------------------------------------------------------
%  Abstract
% -------------------------------------------------------

\begin{center}
\textbf{چکیده}
\end{center}

در تصاویر سی‌تی بدون تزریق ماده‌ی حاجب، تباین میان بافت‌های نرم کم است و عملکرد مدل‌های بخش‌بندی، به‌ویژه در عروق، نسبت به تصاویر دارای ماده‌ی حاجب کاهش می‌یابد. کمبود داده‌ی برچسب‌دار در این مرحله و دشواری هم‌ترازسازی مراحل مختلف پویش، ناشی از تغییرات آناتومیکی و فاصله‌ی زمانی تصویربرداری، این مسئله را پیچیده‌تر می‌کند.

در این پژوهش، ابتدا یک پیش‌پردازش آناتومی‌محور برای هم‌ترازسازی مراحل پویش ارائه شد که شمار خطاهای بزرگ هم‌ترازی (جابه‌جایی بیش از \عدد{۱۰} میلی‌متر در مرکز اندام) را برای دقیق‌ترین روش بررسی‌شده، \lr{DEEDS}، \عدد{۶۵} درصد کاهش داد. جفت‌های هم‌ترازشده برای آموزش مدل پیشنهادی \lr{PLC-Net} به کار رفتند. این مدل تصویر بدون تزریق را همراه با مرحله‌ی هدف پویش و سطح شدت هدف اندام‌ها برحسب واحد هانسفیلد دریافت و تصویر همان بیمار را در مرحله‌ی هدف تولید می‌کند. آزمایش‌ها روی مجموعه‌داده‌ی \lr{VinDr-Multiphase} با تقسیم بیمارمحور \عدد{۹۷}، \عدد{۲۰} و \عدد{۲۰} مورد برای آموزش، اعتبارسنجی و آزمون انجام شد و خروجی‌ها با معیارهای تصویری و بخش‌بندی خودکار به‌عنوان وظیفه‌ی پایین‌دستی ارزیابی شدند.

خروجی \lr{PLC-Net} عملکرد بخش‌بندی در ناحیه‌ی ورید باب را از میانگین \lr{Dice} برابر با \عدد{۰٫۵۵} در تصویر بدون تزریق به \عدد{۰٫۷۲} افزایش داد. در میان معیارهای تصویری، خطای میانه‌ی عدد هانسفیلد در سطح اندام بیشترین هم‌بستگی را با بهبود بخش‌بندی داشت. \lr{PLC-Net} با خطای \عدد{۱۴٫۵۵} \lr{HU} در این معیار، بهترین عملکرد را میان روش‌های مقایسه‌شده به دست آورد؛ در حالی که روش نزدیک، \lr{ResViT}، با حدود \عدد{۹} برابر پارامتر بیشتر به \عدد{۱۵٫۱۷} \lr{HU} رسید. این نتایج نشان می‌دهند که رویکرد پیشنهادی توانسته است با یک مدل شرطی کم‌حجم، تصاویر دارای ماده‌ی حاجب سودمندی برای وظیفه‌ی پایین‌دستی تولید کند و راهی برای استفاده از تصاویر بدون تزریق در بخش‌بندی اندام‌های شکمی بردارد.


\bigskip\noindent\textbf{کلیدواژه‌ها}:
بخش‌بندی چندعضوی، سی‌تی شکمی، تصویر بدون‌تزریق، هم‌ترازسازی بین‌پویشی، هم‌نهادی شرطی تصویر، ارزیابی پایین‌دستی
